\documentclass[12pt,a4paper,brazil]{article}
\usepackage[portuguese]{babel}
\usepackage[T1]{fontenc}
\usepackage[hidelinks]{hyperref}
\usepackage[onehalfspacing]{setspace}
\usepackage{amsmath, amssymb}
\usepackage{fancyhdr}
\usepackage{float}
\usepackage{geometry}
\usepackage{graphicx}
\usepackage{pgfplots}
\usepackage{pgfplotstable}
\usepackage{tikz}
\usepackage{titling}
\usepackage{xparse}
\usepackage{xspace}

\usetikzlibrary{arrows.meta, positioning}
\usetikzlibrary{graphs}
\usetikzlibrary{calc}
\pgfplotsset{compat=1.18}

\setlength{\parindent}{30pt}
\renewcommand{\familydefault}{lmss}
\newcommand{\reffig}[1]{\hyperref[#1]{Figura \ref{#1}}}
\newcommand{\reduct}{\texorpdfstring{$\leq_p$}{≤ₚ}\xspace}
\NewDocumentCommand{\grafico}{m m m O{north west}}{%
% #1 - Nome do arquivo .csv
% #2 - Coluna do eixo y
% #3 - Título do eixo y
% #4 - Posição da legenda

    \begin{tikzpicture}
        \begin{axis}[
            xlabel={Número de vértices},
            ylabel={#3},
            legend pos=#4,
            xtick={5,10,15,20,25},
            width=14cm,
            height=7cm,
            ymode=log,
        ]
            \addplot table[x=n,y=bmk_#2,col sep=comma] {resultados/results_#1.csv};
            \addplot table[x=n,y=sat_#2,col sep=comma] {resultados/results_#1.csv};
            \legend{BMK, SAT}
        \end{axis}
    \end{tikzpicture}
}

\predate{}
\postdate{}
\title{Caminhos Hamiltonianos: Um Estudo Teórico e Experimental}
\author{Izak Alves Gama}
\date{}

\begin{document}
    \maketitle
    
    \begin{abstract}
Este trabalho apresenta um estudo teórico e experimental sobre o Problema do Caminho Hamiltoniano em grafos não direcionados e conexos. Teoricamente, são apresentadas as definições formais do problema, demonstra-se seu pertencimento à classe NP e prova-se sua NP-completude por meio de reduções polinomiais a partir do 3-SAT, além de uma discussão sobre métodos para resolvê-lo. Experimentalmente, são comparadas duas abordagens exatas: uma baseada em programação dinâmica (Bellman--Held--Karp) e outra baseada em redução para SAT com um solucionador Conflict-Driven Clause Learning, avaliadas quanto ao tempo de execução e ao consumo de memória em instâncias de diferentes tamanhos e densidades. Os resultados indicam que a abordagem baseada em programação dinâmica apresenta melhor desempenho temporal para as instâncias consideradas, enquanto a redução para SAT sofre crescimento mais acentuado do tempo de execução; e que ambas apresentam crescimento exponencial no uso de memória, ainda que na redução para SAT seja mais suave.
\end{abstract}
    \section{Introdução}

O Problema do Caminho Hamiltoniano é um dos problemas mais conhecidos da ciência da computação, inclusive sendo um dos problemas NP-Completos de Karp \cite{Karp1972}. Do ponto de vista histórico, o conceito de caminho hamiltoniano remonta ao século XIX, vindo de um jogo criado por William Rowan Hamilton, e desde então tornou-se um objeto de estudo clássico na teoria dos grafos \cite{West2000}. Do ponto de vista da ciência da computação, sua relevância vem de seu valor teórico, como exemplo de problema intratável no pior caso, e também de suas relações com outros problemas de otimização e sequenciamento, como o Problema do Caixeiro Viajante, planejamento de rotas, ordenação de tarefas e análise de redes.

Seja $G = (V, E)$ um grafo não direcionado. Um caminho simples $P \in G$ é um subgrafo definido por $P = (V', E')$, onde $V' = \{v_1, v_2, \ldots,v_n\}$, $E' = \{v_1v_2, v_2v_3, \ldots,$ $v_{n-1}v_n\}$ e, $\forall v_i, v_j \in V'$, $i \neq j \Rightarrow v_i \neq v_j$. De forma mais intuitiva, um caminho simples corresponde a uma trajetória no grafo que se inicia em um vértice, percorre uma sequência de vértices adjacentes e termina em um vértice final, sem repetir vértices ao longo do percurso. Se um grafo $G$ possui um caminho simples $P$ que contém todos os seus vértices, então $P$ é chamado de Caminho Hamiltoniano.

Determinar a existência de um caminho dessa natureza em um grafo qualquer é um problema computacionalmente difícil. Mais precisamente, na sua versão de decisão, dado um grafo não direcionado $G = (V,E)$, deseja-se decidir se existe um caminho hamiltoniano em $G$. Isso define o Problema do Caminho Hamiltoniano (HAMPATH), que é conhecido por ser NP-completo.

Neste trabalho, estuda-se o Problema do Caminho Hamiltoniano de forma teórica e experimental. Teoricamente, são apresentadas as definições formais, a prova de pertencimento do problema à classe NP e sua NP-completude por meio de reduções a partir do 3-SAT para o Problema do Caminho Hamiltoniano Dirigido (DHAMPATH), que por sua vez é reduzido para HAMPATH. Também são discutidos métodos para resolução do problema, desde abordagens exatas até heurísticas e aproximações.

Experimentalmente, são comparadas duas abordagens exatas para a resolução do problema: um algoritmo baseado em programação dinâmica (Bellman-Held-Karp - BMK) e uma abordagem baseada em redução para SAT resolvida por um solucionador CDCL (\textit{Conflict-Driven Clause Learning}). O objetivo é analisar empiricamente o desempenho dessas abordagens em termos de tempo de execução e consumo de memória, discutindo suas vantagens e limitações.

    \section{NP-Completude}

Nesta seção é demonstrado formalmente que o Problema do Caminho Hamiltoniano é NP-completo. Para isso, inicialmente mostra-se que o problema pertence à classe NP. Em seguida, é provado que o problema é NP-Completo por meio de reduções polinomiais, começando com uma redução do 3-SAT para o Problema do Caminho Hamiltoniano Dirigido que, por sua vez, é reduzido para não direcionado.

Antes de tudo, é necessário mostrar que Problema do Caminho Hamiltoniano (HAMPATH) pertence à classe NP. A instância é definida por um grafo não dirigido $G = (V, E)$ e o certificado por um caminho simples $P = (V', E')$, conforme definido anteriormente. A verificação de que $P$ é um caminho hamiltoniano em $G$ pode ser realizada em tempo polinomial da seguinte forma:

\begin{itemize}
    \item Verificar se $V' = V$, isto é, se $P$ e $G$ possuem os mesmos vértices.
    \item Para cada par consecutivo $v_i, v_{i+1}\in V'$, verificar se a aresta $v_iv_{i+1} \in E$.
\end{itemize}

Verificar se $V'=V$  pode ser feito em tempo $O(|V| \log |V| + |V'| \log |V'|)$, bastando ordenar $V$ e $V'$, e verificar se $|V|=|V'|$ e $\forall i= 1,\ldots,|V|$ se $v_i=v'_i$, com $v_i\in V$ e $v_i'\in V'$. Já a verificação da existência das arestas pode ser feita em tempo $O(|V|)$, iterando em $V'$ e assumindo que as arestas foram representadas como uma matiz/lista de adjacência. Logo, é possível verificar em tempo polinomial se um certificado $P$ é um caminho hamiltoniano em $G$ e portanto, HAMPATH pertence à classe NP.

Para provar que HAMPATH é NP-Completo, serão realizadas duas reduções: a primeira é do Problema da 3-Satisfatibilidade (3-SAT) (um problema NP-Completo \cite{Karp1972, cormen2012} bem conhecido) para o Problema do Caminho Hamiltoniano Dirigido (DHAMPATH), que por sua vez será reduzido para HAMPATH. De maneira análoga ao que foi provado anteriormente para HAMPATH, pode-se provar que DHAMPATH também pertence à classe NP.

\subsection{3-SAT \reduct DHAMPATH \texorpdfstring{\cite{Kleinberg2006}}{(Kleinberg; Tardos, 2006)}}

Seja uma instância arbitrária do 3-SAT com variáveis $x_1,\ldots,x_n$ e cláusulas $C_1,\ldots,C_k$. A fim de auxiliar a compreensão, será tomada como exemplo a seguinte expressão $E$ como instância do 3-SAT a ser reduzida: $$E=(x\lor\lnot y\lor z)\land(\lnot x\lor w\lor\lnot z)$$

A ideia é, a princípio, construir um grafo que contenha $2^n$ caminhos hamiltonianos possíveis, onde cada um deles representam as $2^n$ atribuições possíveis de valores às variáveis. Constrói-se $n$ caminhos $P_1, \ldots, P_n$ para cada variável, onde cada $P_i$ possui vértices $v_{i1}, v_{i2}, \ldots,v_{ib}$, onde $b>k$.

Uma variável pode estar em todas as $k$ cláusulas, então vamos assumir que $b=2k+2$, isto é, todo $P_i$ terá 2 vértices para cada cláusula e mais 2 extras (isso fará sentido adiante). Para cada vértice $v_{ij}$, com $1 \le j < b$, existem as arestas $v_{ij}v_{i,j+1}$ e $v_{i,j+1}v_{ij}$, ou seja, cada caminho $P_i$ pode ser percorrida da esquerda para a direita e vice-versa.

Os caminhos são conectados da seguinte forma: para cada $i=1,\ldots, n-1$, são criadas as arestas $v_{i1}v_{i+1,1}$, $v_{i1}v_{i+1,b}$, $v_{ib}v_{i+1,1}$ e $v_{ib}v_{i+1,b}$. São adicionados ao grafo 2 novos vértices $s$ e $t$ (para garantir que qualquer caminho hamiltoniano comece em $s$ e termine em $t$), e definidos as arestas $sv_{11}$, $sv_{1b}$, $tv_{n1}$ e $tv_{nb}$. Até este ponto, a estrutura do grafo está ilustrada abaixo na \reffig{fig:3-sat-para-chd_parcial}.

\begin{figure}[H]
    \centering
    \input{figuras/3-sat-para-chd_parcial}
    \caption{Grafo parcial formado da redução de 3-SAT para DHAMPATH}
    \label{fig:3-sat-para-chd_parcial}
\end{figure}

Seguindo o processo descrito anteriormente para $E$, com $n=4$ e $k=2$, são construídos 4 caminhos com $b=2\cdot2+2=6$ vértices cada. Adicionando os vértices $s$ e $t$, e conectando todos os vértices como já descrito, o grafo parcial para $E$ está ilustrado abaixo na \reffig{fig:3-sat-para-chd_parcial_exemplo}.

\begin{figure}[H]
    \centering
    \begin{tikzpicture}[
    node/.style={circle, draw, minimum size=6mm},
    >=Stealth
]

\def\xsep{3}
\def\ysep{1.75}

% ---------- Colunas ------------
\foreach \i [evaluate=\i as \yy using {-(\i-1)*\ysep}] in {1,...,6} {
  \node[node] (px\i) at (0,\yy) {$v_{x,\i}$};
  \node[node] (py\i) at (\xsep,\yy) {$v_{y,\i}$};
  \node[node] (pz\i) at (2*\xsep,\yy) {$v_{z,\i}$};
  \node[node] (pw\i) at (3*\xsep,\yy) {$v_{w,\i}$};
}

\foreach \i in {x,y,z,w} {
    \node at ($(p\i6)+(0,-1.2)$) {$P_\i$};
}

% ----------- Arestas internas -----------
\foreach \col in {px,py,pz,pw}{
  \foreach \i in {1,...,5}{
    \def\iinc{\the\numexpr\i+1\relax}
    \draw[->,bend left=15] (\col\i) to (\col\iinc);
    \draw[->,bend left=15] (\col\iinc) to (\col\i);
  }
}

% Vértices s e t
\path let \p1 = (px1), \p2 = (px6) in
  node[node] (s) at (-\xsep,{(\y1+\y2)/2}) {$s$};

\path let \p1 = (pw1), \p2 = (pw6) in
  node[node] (t) at (4*\xsep,{(\y1+\y2)/2}) {$t$};

\draw[->] (s) -- (px1);
\draw[->] (s) -- (px6);
\draw[->] (pw1) -- (t);
\draw[->] (pw6) -- (t);

% Conexões entre as colunas

\foreach \a/\b in {x/y,y/z,z/w} {
    \draw[->] (p\a1) -- (p\b1);
    \draw[->] (p\a1) -- (p\b6);
    \draw[->] (p\a6) -- (p\b1);
    \draw[->] (p\a6) -- (p\b6);
}

\end{tikzpicture}

    \caption{Grafo parcial formado da redução da instância $E$ do 3-SAT para DHAMPATH}
    \label{fig:3-sat-para-chd_parcial_exemplo}
\end{figure}

A ideia é que cada caminho hamiltoniano seja identificado da seguinte maneira: se o caminho percorrer $P_i$ da esquerda para direita, então $x_i$ é verdadeiro; caso contrário, é falso. Para modelar isso, serão adicionados novos $k$ vértices (um para cada cláusula) e a instância do 3-SAT será satisfatível se, e somente se, restar algum caminho hamiltoniano no grafo.

Seja $c_j$ o vértice da cláusula $C_j$. Para cada termo $t\in C_j$, se $t=x_i$, então adicionaremos as arestas $v_{i,2j}c_{j}$ e $c_{j}v_{i,2j+1}$; se $t=\lnot x_i$, então adicionaremos as arestas $v_{i,2j+1}c_{j}$ e $c_{j}v_{i,2j}$. A \reffig{fig:3-sat-para-chd_total} mostra visualmente como $c_j$ é adicionado ao grafo.

\begin{figure}[H]
    \centering
    \begin{tikzpicture}[
    node/.style={circle, draw, minimum size=6mm},
    dot/.style={inner sep=0pt},
    >=Stealth
]

    \def\xsep{3}
    \def\ysep{1.8}
    \def\xn{2.6*\xsep}
    \def\ycenter{-1.85*\ysep}
    
    % ---------- vértice s ----------
    \node[node] (s) at (-\xsep,\ycenter) {$s$};
    
    % ---------- Coluna P1 ----------
    \node[node] (p11) at (0,0) {$v_{1,1}$};
    \node[node] (p12) at (0,-\ysep) {$v_{1,2}$};
    \node[node] (p13) at (0,-2*\ysep) {$v_{1,3}$};
    \node at (0,-2.85*\ysep) {$\vdots$};
    \node[node] (p1k) at (0,-3.7*\ysep) {$v_{1,k}$};
    
    \node at ($(p1k)+(0,-1.2)$) {$P_1$};
    
    % ---------- Coluna P2 ----------
    \node[node] (p21) at (\xsep,0) {$v_{2,1}$};
    \node[node] (p22) at (\xsep,-\ysep) {$v_{2,2}$};
    \node[node] (p23) at (\xsep,-2*\ysep) {$v_{2,3}$};
    \node at (\xsep,-2.85*\ysep) {$\vdots$};
    \node[node] (p2k) at (\xsep,-3.7*\ysep) {$v_{2,k}$};
    
    \node at ($(p2k)+(0,-1.2)$) {$P_2$};
    
    % ---------- Reticências horizontais ----------
    \node (hdots) at ({(\xsep+\xn)/2}, -2*\ysep) {$\cdots$};
    
    % Pontos auxiliares para as arestas
    \node[dot] (L) at ($(hdots)+(-0.6,0)$) {};
    \node[dot] (R) at ($(hdots)+(0.6,0)$) {};
    
    % ---------- Coluna Pn ----------
    \node[node] (pn1) at (\xn,0) {$v_{n,1}$};
    \node[node] (pn2) at (\xn,-\ysep) {$v_{n,2}$};
    \node[node] (pn3) at (\xn,-2*\ysep) {$v_{n,3}$};
    \node at (\xn,-2.85*\ysep) {$\vdots$};
    \node[node] (pnk) at (\xn,-3.7*\ysep) {$v_{n,k}$};
    
    \node at ($(pnk)+(0,-1.2)$) {$P_n$};
    
    % ---------- vértice t ----------
    \node[node] (t) at (\xn+\xsep,\ycenter) {$t$};
    
    % ---------- Arestas internas ----------
    \foreach \a/\b in {p11/p12,p12/p13,p21/p22,p22/p23,pn1/pn2,pn2/pn3}{
      \draw[->,bend left=20] (\a) to (\b);
      \draw[->,bend left=20] (\b) to (\a);
    }
    
    % ---------- Arestas de P1 para P2 ----------
    \draw[->] (p11) -- (p21);
    \draw[->] (p11) -- (p2k);
    \draw[->] (p1k) -- (p21);
    \draw[->] (p1k) -- (p2k);
    
    % ---------- Adicionado manualmente ---------
    \def\edx{1.2}
    \def\eix{0.25}
    \def\eiy{0.55*\ysep}
    
    \draw[-] (p21) -- ++(\edx, 0);
    \draw[-] (p21) -- ++(\edx, -\ysep);
    \draw[-] (p2k) -- ++(\edx, 0);
    \draw[-] (p2k) -- ++(\edx, \ysep);
    
    \draw[<-] (pn1) -- ++(-\edx, 0);
    \draw[<-] (pn1) -- ++(-\edx, -\ysep);
    \draw[<-] (pnk) -- ++(-\edx, 0);
    \draw[<-] (pnk) -- ++(-\edx, \ysep);
    
    \node (hdotsTop) at ($(p21)!0.5!(pn1)$) {$\cdots$};
    \node (hdotsBottom) at ($(p2k)!0.5!(pnk)$) {$\cdots$};
    
    \foreach \a in {p13,p23,pn3}{
      \draw[-,bend left=10] (\a) to ++(\eix, -\eiy);
      \draw[<-,bend right=10] (\a) to ++(-\eix, -\eiy);
    }
    
    \foreach \a in {p1k,p2k,pnk}{
      \draw[<-,bend right=10] (\a) to ++(\eix, \eiy);
      \draw[-,bend left=10] (\a) to ++(-\eix, \eiy);
    }
    
    % ---------- Arestas com s e t ----------
    \draw[->] (s) -- (p11);
    \draw[->] (s) -- (p1k);
    \draw[->] (pn1) -- (t);
    \draw[->] (pnk) -- (t);
    
    % ------------- Vértice c1 ---------------
    \node[node] (c1) at ({0.5*\xsep}, {1.25*\ysep}) {$c_1$};
    
    \draw[->,dashed,out=45,in=225] (p12) to (c1);
    \draw[->,dashed,out=250,in=45] (c1) to (p13);
    
    \draw[->,dashed,out=300,in=135] (c1) to (p22);
    \draw[->,dashed,out=135,in=270] (p23) to (c1);
    
    \draw[->,dashed] (pn2) -- (c1);
    \draw[->,dashed] (c1) -- (pn3);


\end{tikzpicture}

    \caption{Grafo formado da redução de 3-SAT para DHAMPATH. Apenas o $c_1$ foi adicionado para uma melhor visualização, e as arestas que o ligam foram adicionadas sem perca de generalidade.}
    \label{fig:3-sat-para-chd_total}
\end{figure}

Tomando como exemplo a primeira cláusula de $E$ ($x\lor\lnot y\lor z$), o caminho hamiltoniano deve percorrer $P_x$ da esquerda para a direita, ou deve percorrer $P_y$ da direita para a esquerda, ou deve percorrer $P_z$ da esquerda para a direita. Então adiciona-se o vértice $c_1$ e as arestas $v_{x,2}c_1$, $c_1v_{x,3}$, $v_{y,3}c_1$, $c_1v_{y,2}$, $v_{z,2}c_1$, $c_1v_{z,3}$. De maneira análoga adiciona-se $c_2$ para a segunda cláusula de $E$ e suas respectivas arestas. A ilustração do grafo formado pela redução de $E$ para DHAMPATH está abaixo na \reffig{fig:3-sat-para-chd_total_exemplo} 

\begin{figure}[H]
    \centering
    \begin{tikzpicture}[
    node/.style={circle, draw, minimum size=6mm},
    >=Stealth
]

\def\xsep{3}
\def\ysep{1.75}

% ---------- Colunas ------------
\foreach \i [evaluate=\i as \yy using {-(\i-1)*\ysep}] in {1,...,6} {
  \node[node] (px\i) at (0,\yy) {$v_{x,\i}$};
  \node[node] (py\i) at (\xsep,\yy) {$v_{y,\i}$};
  \node[node] (pz\i) at (2*\xsep,\yy) {$v_{z,\i}$};
  \node[node] (pw\i) at (3*\xsep,\yy) {$v_{w,\i}$};
}

\foreach \i in {x,y,z,w} {
    \node at ($(p\i6)+(0,-1.2)$) {$P_\i$};
}

% ----------- c1 e c2 -----------
\node[node] (c1) at (0.5*\xsep,1.25*\ysep) {$c_1$};
\node[node] (c2) at (2.5*\xsep,1.25*\ysep) {$c_2$};

% ----------- Arestas internas -----------
\foreach \col in {px,py,pz,pw}{
  \foreach \i in {1,...,5}{
    \def\iinc{\the\numexpr\i+1\relax}
    \draw[->,bend left=15, red] (\col\i) to (\col\iinc);
    \draw[->,bend left=15] (\col\iinc) to (\col\i);
  }
}

\draw[->,bend left=15] (px2) to (px3);
\draw[->,bend left=15] (pw4) to (pw5);

% ----------- Conexões com c1 -----------
% Px
\draw[->,dashed,out=45,in=210,red] (px2) to (c1);
\draw[->,dashed,out=240,in=45,red] (c1) to (px3);

% Py (invertido)
\draw[->,dashed,out=300,in=135] (c1) to (py2);
\draw[->,dashed,out=135,in=270] (py3) to (c1);

% Pz
\draw[->,dashed] (pz2) -- (c1);
\draw[->,dashed] (c1) -- (pz3);

% ----------- Conexões com c2 -----------
% Px (invertido)
\draw[->,dashed,out=180,in=30] (c2) to (px4);
\draw[->,dashed,out=30,in=210] (px5) to (c2);

% Pz (invertido)
\draw[->,dashed,out=240,in=45] (c2) to (pz4);
\draw[->,dashed,out=30,in=270] (pz5) to (c2);

% Pw (normal)
\draw[->,dashed,out=30,in=0,red] (pw4) to (c2);
\draw[->,dashed,out=330,in=30,red] (c2) to (pw5);

% Vértices s e t
\path let \p1 = (px1), \p2 = (px6) in
  node[node] (s) at (-\xsep,{(\y1+\y2)/2}) {$s$};

\path let \p1 = (pw1), \p2 = (pw6) in
  node[node] (t) at (4*\xsep,{(\y1+\y2)/2}) {$t$};

\draw[->,red] (s) -- (px1);
\draw[->] (s) -- (px6);
\draw[->] (pw1) -- (t);
\draw[->,red] (pw6) -- (t);

% Conexões entre as colunas

\foreach \a/\b in {x/y,y/z,z/w} {
    \draw[->] (p\a1) -- (p\b1);
    \draw[->] (p\a1) -- (p\b6);
    \draw[->, red] (p\a6) -- (p\b1);
    \draw[->] (p\a6) -- (p\b6);
}

\end{tikzpicture}

    \caption{Grafo formado da redução da instância $E$ do 3-SAT para DHAMPATH. As arestas que formam um caminho hamiltoniano no grafo está destacado em vermelho, indicando que $E$ é satisfatível.}
    \label{fig:3-sat-para-chd_total_exemplo}
\end{figure}

Isso completa a construção do grafo e, consequentemente, a redução. O grafo possui $n$ caminhos $P_i$ de $b=2k+2$ vértices, $k$ vértices de cláusulas, um de início e outro final, totalizando $2nk+2n+k+2 = \Theta(nk)$ vértices. Em cada caminho $P_i$, há 2 arestas entre dois vértices adjacentes (ida e volta), e mais 2 em $v_{i1}$ e $v_{i,b}$ para conexão com $P_{i+1}$ (com exceção de $P_n$), além de mais 2 conectando $s$ com $P_1$ e 2 conectando $P_n$ com $t$, e mais 6 para cada $c_j$, totalizando $4nk+6n+6k=\Theta(nk)$ arestas. Considerando que a criação dos vértices quanto das arestas são realizadas em tempo $\Theta(nk)$ para ambos, pode-se afirmar que a redução é polinomial.

Agora será provado que uma instância do 3-SAT é satisfatível se, e somente se, o grafo construído a partir da redução possuir um caminho hamiltoniano. Dado uma instância $I=(\{x_1,\ldots,x_n\},;\{C_1,\ldots,C_k\})$ qualquer do 3-SAT, supõe-se que exista uma atribuição das variáveis que a satisfaça. Se $x_i$ for verdadeiro, espera-se que $P_i$ seja ``percorrido'' da esquerda para a direita; caso contrário, da direita para a esquerda. Para cada $C_j$, como ela é satisfeita pela atribuição, haverá ao menos um caminho $P_i$ que será percorrido na direção ``correta'' em relação à $c_j$, e podemos inseri-lo no percurso por meio das arestas incidentes em $v_{i,2j}$ e $v_{i,2j+1}$.

Por outro lado, supõe-se que exista um caminho hamiltoniano $P$ no grafo. Se $P$ entra em um vértice $c_j$ por uma aresta proveniente de $v_{i,2j}$, ele deve sair por uma aresta para $v_{i,2j+1}$. Caso contrário, $v_{i,2j+1}$ terá apenas um vizinho não visitado restante, ou seja, $v_{i,2j+2}$ e, portanto, o percurso não poderá visitar esse vértice e ainda manter a propriedade hamiltoniana. Simetricamente, se entra por $v_{i,2j+1}$ deve-se sair por $v_{i,2j}$.

Os vértices imediatamente anteriores e posteriores a $c_j$ no caminho $P$ são unidos por uma aresta $e$ no grafo; portanto, se $c_j$ for removido do caminho e for inserido $e$ cada $j$, obtém-se um caminho hamiltoniano $P'$ no subgrafo formado ao remover $c_1,\ldots,c_k$ (grafo original, antes dos vértices de cláusula serem adicionados). Qualquer caminho hamiltoniano neste subgrafo deve percorrer cada $P_i$ completamente em uma direção ou outra.

Assim, $P'$ é usado para definir a seguinte atribuição de verdade para a instância do 3-SAT: Se $P'$ percorrer $P_i$ da esquerda para a direita, $x_i = 1$; caso contrário, $xi = 0$. Como o caminho maior $P$ foi capaz de visitar cada vértice de cláusula $c_j$, pelo menos um dos caminhos foi percorrido na direção “correta” em
relação ao vértice $c_j$, e, portanto, a atribuição definida até aqui satisfaz todas as cláusulas. 

Tendo estabelecido que a instância do 3-SAT é satisfatível se, e somente se, o grafo possuir um caminho hamiltoniano, portanto, DHAMPATH é NP-Completo.


\subsection{DHAMPATH \reduct HAMPATH}

Seja $G=(V,E)$ um grafo dirigido qualquer. A fim de auxiliar a compreensão, será tomado como exemplo o seguinte grafo como instância do DHAMPATH a ser reduzida para HAMPATH:

\begin{figure}[H]
    \centering
    \begin{tikzpicture}[
    vertex/.style={circle, draw, minimum size=9mm, inner sep=0pt},
    edge/.style={->, thick},
    pathedge/.style={->, very thick, red}
]

    \node[vertex] (a) at (0,1.5) {a};
    \node[vertex] (b) at (2,2.5) {b};
    \node[vertex] (c) at (2,0.5) {c};
    \node[vertex] (d) at (4,1.5) {d};
    
    % Arestas do grafo
    \draw[edge] (a) -- (b);
    \draw[edge] (a) -- (c);
    \draw[edge] (b) -- (c);
    \draw[edge] (b) -- (d);
    \draw[edge] (c) -- (d);
    
    % Caminho destacado a -> b -> c -> d
    \draw[pathedge] (a) -- (b);
    \draw[pathedge] (b) -- (c);
    \draw[pathedge] (c) -- (d);

\end{tikzpicture}
    \caption{Grafo dirigido a ser reduzido para não dirigido. As arestas em vermelho destacam o caminho hamiltoniano do grafo.}
    \label{fig:chd-para-chnd_grafo_exemplo}
\end{figure}

\noindent A ideia é transformar cada vértice de G em uma pequena subestrutura que force o caminho a ter uma direção lógica, ainda que as arestas não tenham direção no novo grafo $G'$.

Para cada vértice $v\in V$, são criados 3 vértices em $G'$: $v_{in}$ (entrada), $v_{mid}$ (meio) e $v_{out}$ (saída). Depois são adicionadas as seguintes arestas em $G'$: $v_{in}v_{mid}$ e $v_{mid}v_{out}$. Após isso, para cada vértice $uv\in E$ é adicionado a aresta $u_{out}v_{in}$. A \reffig{fig:chd-para-chnd_grafo_reduzido} ilustra como o grafo de exemplo ficou após esses processos.

\begin{figure}[H]
    \centering
    \begin{tikzpicture}[
    vertex/.style={circle, draw, minimum size=9mm, inner sep=0pt},
    edge/.style={-, thick},
    pathedge/.style={-, very thick, red}
]

    % Posições globais do diamante
    \node (A) at (-2,1.5) {};
    \node (B) at (2,2.75) {};
    \node (C) at (2,0.25) {};
    \node (D) at (6,1.5) {};
    
    % Deslocamento interno maior
    \def\d{1.5}
    
    % a_in, a_mid, a_out
    \node[vertex] (ain)  at ($(A)+(-\d,0)$) {$a_{in}$};
    \node[vertex] (amid) at (A)            {$a_{mid}$};
    \node[vertex] (aout) at ($(A)+(\d,0)$)  {$a_{out}$};
    
    % b_in, b_mid, b_out
    \node[vertex] (bin)  at ($(B)+(-\d,0)$) {$b_{in}$};
    \node[vertex] (bmid) at (B)            {$b_{mid}$};
    \node[vertex] (bout) at ($(B)+(\d,0)$)  {$b_{out}$};
    
    % c_in, c_mid, c_out
    \node[vertex] (cin)  at ($(C)+(-\d,0)$) {$c_{in}$};
    \node[vertex] (cmid) at (C)            {$c_{mid}$};
    \node[vertex] (cout) at ($(C)+(\d,0)$)  {$c_{out}$};
    
    % d_in, d_mid, d_out
    \node[vertex] (din)  at ($(D)+(-\d,0)$) {$d_{in}$};
    \node[vertex] (dmid) at (D)            {$d_{mid}$};
    \node[vertex] (dout) at ($(D)+(\d,0)$)  {$d_{out}$};
    
    % Arestas internas
    \draw[edge] (ain) -- (amid) -- (aout);
    \draw[edge] (bin) -- (bmid) -- (bout);
    \draw[edge] (cin) -- (cmid) -- (cout);
    \draw[edge] (din) -- (dmid) -- (dout);
    
    % Arestas conforme o grafo original
    \draw[edge] (aout) -- (bin);
    \draw[edge] (aout) -- (cin);
    \draw[edge] (bout) -- (cin);
    \draw[edge] (bout) -- (din);
    \draw[edge] (cout) -- (din);
    
    % Caminho destacado
    \draw[pathedge] (ain) -- (amid) -- (aout) -- (bin) -- (bmid) -- (bout)
                   -- (cin) -- (cmid) -- (cout) -- (din) -- (dmid) -- (dout);

\end{tikzpicture}

    \caption{Grafo dirigido reduzido para não dirigido. As arestas em vermelho destacam o caminho hamiltoniano do grafo.}
    \label{fig:chd-para-chnd_grafo_reduzido}
\end{figure}

Com isso, a redução do DHAMPATH para HAMPATH está concluída. Em $G$ O número de vértices é $3|V|=\Theta(|V|)$, enquanto o número de arestas é $2|V|+|E|=\Theta(|V|+|E|)$ (2 arestas internas por vértice mais as arestas originais). Considerando que a construção dos vértices e arestas em $G'$ são realizadas em tempo $\Theta(|V|)+\Theta(|V|+|E|)=\Theta(|V|+|E|)$, pode-se afirmar que a redução é polinomial.

Agora será provado que $G$ possui um caminho hamiltoniano se, e somente se, $G'$ possui um caminho hamiltoniano. Supondo que exista um caminho hamiltoniano em $G$, ele pode ser organizado em uma sequência de vértices $v_1,v_2,\ldots,v_n$, onde $\forall i=1,\ldots,n-1,\;v_iv_{i+1}\in E$. A partir dessa sequência, pode-se construir a sequência $v_{1}^{in},v_{1}^{mid},v_{1}^{out},v_{2}^{in},v_{2}^{mid},v_{2}^{out},\ldots,v_{n}^{in},v_{n}^{mid},v_{n}^{out}$, onde cada 2 vértices adjacentes são ligados por uma aresta. As arestas $v_{i}^{in}v_{i}^{mid}$ e $v_{i}^{mid}v_{i}^{out}$ existem por construção, enquanto $v_{i}^{out},v_{i+1}^{in}$ existe porque $v_iv_{i+1}\in E$. Logo, Esse caminho cobre todos os $3n$ vértices uma única vez.

Por outro lado, supondo que exista um caminho hamiltoniano em $G'$, em cada vértice $v\in V$ o seu representante $v_{mid}$ está conectado apenas à $v_{in}$ e $v_{out}$; logo, qualquer caminho hamiltoniano que visite $v_{mid}$ deve entrar por um vizinho e sair imediatamente pelo outro, ou seja, $v_{in},v_{mid}$ e $v_{out}$ devem ser percorridos consecutivamente. A única forma de conectar um bloco $\{u_{in},u_{mid},u_{out}\}$ a outro bloco $\{v_{in},v_{mid},v_{out}\}$ é através de uma aresta $u_{out}v_{in}$ ou $v_{in},u_{out}$, que por sua vez, por construção, só existem se houver uma aresta $uv$ ou $vu$, respectivamente, em $G$. Portanto, a sequência de blocos visitados em $G'$ corresponde diretamente à sequência de vértices em $G$.

Tendo estabelecido que a instância do DHAMPATH possui um caminho hamiltoniano se, e somente se, a sua redução para HAMPATH também possuir um caminho hamiltoniano, portanto, HAMPATH é NP-Completo.
    \section{Métodos de resolução} \label{sec:metodos-resolucao}

Nesta seção são apresentados e discutidos métodos para a resolução do HAMPATH. São abordados métodos exatos (corretos, mas computacionalmente caro), heurísticos e algoritmos aproximados (esses dois últimos são rápidos, mas não garantem corretude).

Para resolver o HAMPATH, há diversos métodos baseados nos mais variados paradigmas, como busca exaustiva (\textit{backtracking} e \textit{branch-and-bound}), programação dinâmica (algoritmo de Bellman-Held-Karp (BHK)) e redução para outro problema conhecido (redução para SAT). Para os experimentos deste trabalho (seções \ref{sec:experimentos} e \ref{sec:resultados}) serão utilizados apenas o BHK e redução para SAT em conjunto com DPLL (um algoritmo clássico para resolver o SAT), pois estes métodos possuem complexidades temporais exponenciais, enquanto os métodos baseados em busca exaustiva possuem complexidade fatorial.

Devido à NP-Completude do HAMPATH, também  há várias heurísticas e abordagens aproximadas para tratar o problema com instâncias maiores, onde os métodos exatos seriam computacionalmente muito custosos ou inviáveis. Métodos heurísticos possuem complexidades temporais mais baixas (geralmente polinomiais), mas formalmente não garantem sucesso; enquanto algoritmos aproximados garantem proximalidade da solução ótima, ainda que foquem mais em variantes do problema (por exemplo, encontrar o caminho hamiltoniano mais curto).


\subsection{\textit{Backtracking} e \textit{branch-and-bound}} \label{subsec:backtracking_e_bb}

A abordagem por \textit{backtracking} faz uma pesquisa exaustiva em profundidade que constrói incrementalmente o caminho. Começa num vértice inicial e tenta estender o caminho escolhendo sucessivamente vértices adjacentes ainda não visitados; se chega a um vértice onde não há extensão possível, volta (``\textit{backtrack}'') para tentar outro vértice no passo anterior \cite{geeksforgeeks2025}. Esse algoritmo garante encontrar um caminho hamiltoniano caso exista, mas seu tempo de execução é fatorial \cite{geeksforgeeks2025}, com valores de $n>20$ sendo proibitivos em relação ao tempo de execução.

Já a abordagem por \textit{branch-and-bound} é uma melhoria do \textit{backtracking}, que tenta podar partes do espaço de busca. Nesse contexto, pode-se usar limites (por exemplo, um cálculo de custo mínimo restante) para descartar ramificações que não podem levar a uma solução melhor que a já encontrada \cite{Sergeenko2020}. Na teoria, essa poda pode melhorar o desempenho em casos favoráveis, ainda que possua complexidade fatorial. Entretanto, mesmo com poda, a complexidade ainda cresce muito rápido conforme $n$ aumenta \cite{Sergeenko2020}.


\subsection{Programação dinâmica (BHK) \texorpdfstring{\cite{Held1962,Bellman1962}}{(Held e Karp, 1962; Bellman 1962)}} \label{subsec:programacao_dinamica}

A ideia do BHK é usar um vetor de estado indexado por $(S,v)$, onde $S\subseteq V$ é um conjunto de vértices visitados e $v\in S$ é o último vértice do caminho atual. Define-se $dp[S][v]=1$ se existe um caminho que percorre exatamente os vértices de $S$ e termina em $v$. Inicialmente, para cada vértice $i$, $dp[i][i]=1$ (caminho trivial). A recursão é: para todo $S$ e $v\in S$, $dp[S][v]=\bigvee_{u\in S\backslash\{v\}}(dp[S\backslash\{v\}][u]\land(u,v)\in E)$. Ao final, se existe $v$ tal que $dp[V][v]=1$, então há caminho Hamiltoniano que termina em $v$. Esse algoritmo possui complexidade temporal de $O(n^2 2^n)$ e espacial de $O(n2^n)$.

Embora exponencial, este método é muito mais rápido que os métodos em \ref{subsec:backtracking_e_bb} e consegue resolver instâncias de tamanho maior de forma exata. Suas vantagens são a otimização de repetições de subproblemas e uso eficiente de memória para compartilhar cálculos, enquanto sua limitação ainda é o crescimento exponencial, onde para $n>30$ o tempo e memória tornam-se proibitivos.


\subsection{Redução para SAT}

Outra abordagem é reduzir o caminho Hamiltoniano a um problema de SAT em CNF e usar um solucionador SAT (neste trabalho, será utilizado o CDCL (\textit{Conflict-Driven Clause Learning})). Uma modelagem padrão é introduzir $n^2$ variáveis booleanas $x_{i,j}$ (com $1\le i,j\le n$) onde $x_{i,j}$ indica que o vértice $j$ está na i-ésima posição do caminho \cite{Lyuu2022}. Então constrói-se as seguintes cláusulas em CNF:

\begin{itemize}
    \item $(x_{1j}\lor x_{2j}\lor\ldots\lor x_{nj})$ para cada $j$ (todos os vértices devem aparecer no caminho);
    \item $(\lnot x_{ij}\lor\lnot x_{kj})$ para todo $i,j,k$ com $i\ne k$ (nenhum vértice $j$ deve aparecer 2 vezes no caminho);
    \item $(x_{i1}\lor x_{i2}\lor\ldots\lor x_{in})$ para cada $i$ (cada posição $i$ no caminho deve ser ocupada);
    \item $(\lnot x_{ij}\lor\lnot x_{ik})$ para todo $i,j,k$ com $j\ne k$ (nenhum par de vértices $j$ e $k$ deve ocupar a mesma posição no caminho);
    \item $(\lnot x_{ki}\lor\lnot x_{k+1,j})$ para todo $(i,j)\not\in G$ e $k=1,2,\ldots,n-1$ (vértices $i$ e $j$ não adjacentes não podem ser adjacentes no caminho)
\end{itemize}

\noindent Essas cláusulas asseguram que qualquer atribuição verdadeira corresponde a um caminho válido (e vice-versa) \cite{Lyuu2022}. A construção resulta em $O(n^3)$ cláusulas \cite{Lyuu2022}.

Com a fórmula CNF em mãos, aplica-se o algoritmo CDCL \cite{MarquesSilva1996,Bayardo1997} para solucioná-la. O CDCL é um método clássico eficaz no qual muitos dos solucionadores de SAT do estado da arte se baseiam. É inspirado no DPLL (Davis-Putnam-Logemann-Loveland), mas incorpora técnicas avançadas para aprender com os erros (conflitos) encontrados durante a busca \cite{MarquesSilva2009}.

Como descrito por Marques-Silva, Lynce e Malik (2009) \cite{MarquesSilva2009}, CDCL inicia com um processo cíclico de busca e propagação. A cada etapa, o solucionador escolhe uma variável de decisão e atribui a ela um valor lógico, seguido pela aplicação iterada da regra da cláusula unitária (propagação unitária) para identificar outras variáveis que devem assumir valores específicos como consequência dessa decisão. Se durante esse processo todas os literais de uma cláusula forem avaliados como falsos, a cláusula torna-se insatisfeita; o algoritmo então declara uma situação de conflito e interrompe o avanço da busca para tratar o problema identificado.

Ao identificar um conflito, o CDCL desfaz a última ação e invoca um procedimento de análise de conflito. O algoritmo examina a estrutura das atribuições recentes, geralmente representadas por um grafo de implicação, para diagnosticar a causa raiz da falha lógica. O objetivo é aprender uma nova cláusula através de uma sequência de operações de resolução, adicionando-a à fórmula original; essa nova cláusula serve para impedir que a mesma combinação de atribuições que gerou o erro ocorra novamente no futuro.

Após o aprendizado da cláusula, o algoritmo executa o passo de \textit{backtracking} de maneira estratégica. O CDCL utiliza a cláusula aprendida para realizar um \textit{backtracking} não-cronológico, permitindo que o solucionador pule vários níveis da árvore de decisão para trás de uma só vez, retornando diretamente ao nível de decisão calculado pela análise de conflito (geralmente onde a nova cláusula se torna unitária), o que torna a exploração do espaço de busca significativamente mais eficiente.

Assim como em \ref{subsec:programacao_dinamica}, esse método é mais rápido  do que os métodos em \ref{subsec:backtracking_e_bb}. Além disso, também há a vantagem de aproveitar o estado da arte em solucionadores de SAT (ainda que neste trabalho seja usado um método clássico), pois em vários casos esses solucionadores conseguem desempenho melhor ou igual à busca exaustiva. Porém, a desvantagem ainda é a sua complexidade temporal, que é exponencial ($O(2^n)$).

\subsection{Métodos heurísticos}

Heurísticas buscam caminhos Hamiltonianos de forma mais rápida do que os métodos anteriores, mas sem garantias formais de sucesso. Em geral produzem boas soluções mas podem falhar mesmo quando a solução existe.

Uma estratégia comum é a busca gulosa, onde se inicia-se em um vértice qualquer e estende o caminho sempre que possível, escolhendo algum vizinho não visitado por meio de algum critério (por exemplo, escolhendo vértices de menor grau). Se o algoritmo chega em um ponto onde não é mais possível continuar (``becos sem saída''), termina a tentativa (e eventualmente reiniciar de outro vértice ou com outra ordem) ou tentar heurísticas locais.

Uma outra heurística é a de Pósa \cite{shields2004,POSA1976}, que tenta escapar de ``becos'' usando rotações de caminho, usando o algoritmo PosaSearch. Outras estratégias algoritmos de busca local (tentar melhorar caminhos com trocas de vértices), meta-heurísticas (colônia de formigas, algoritmos genéticos, GRASP etc.) e técnicas específicas de instância \cite{shields2004}.


\subsection{Métodos aproximados}

Como já explicitado antes, algoritmos aproximados focam em variantes de otimização do problema e vêm com garantias de proximidade do ótimo. Monnot \cite{MONNOT2005} trata precisamente de aproximações para o problema do caminho Hamiltoniano de peso máximo com um ou dois vértices fixos ($HPP_s$ e $HPP_{s,t}$, respectivamente). Já o algoritmo de Christofides \cite{Christofides2022} fornece uma garantia de aproximação de $3/2$ para o TSP Métrico usando árvores geradoras mínimas (esse método pode ser adaptado para o HAMPATH).

Forlizzi \textit{et al.} \cite{Forlizzi2005} mostram que, ao combinar uma árvore geradora mínima com um emparelhamento de caminhos cuidadosamente construído, é possível formar um multigrafo com exatamente dois vértices de grau ímpar, correspondentes aos extremos do caminho desejado. Já Wu \textit{et al.} \cite{WU2023} apresentam um algoritmo de aproximação $3/2$ para encontrar um conjunto de $k$ caminhos com o menor peso total nas arestas em grafos completos não direcionados e sem vértices inicial e final prefixados. 
    \section{Experimentos} \label{sec:experimentos}

Esta seção descreve o desenho experimental adotado para avaliar os métodos considerados para o Problema do Caminho Hamiltoniano em grafos não direcionados e conexos. São apresentados os algoritmos comparados, a geração das instâncias de teste, as métricas utilizadas e o ambiente de execução.

Como já introduzido na seção \ref{sec:metodos-resolucao}, neste trabalho foram comparadas duas abordagens: uma baseada no algoritmo de Bellman-Held-Karp (BHK) (programação dinâmica) e redução para SAT (RedSAT), utilizando o \textit{Conflict-Driven Clause Learning} (CDCL) como algoritmo solucionador~\footnote{Foi utilizada a seguinte implementação do CDCL como solucionador: \url{https://github.com/sukrutrao/SAT-Solver-CDCL.git}}. Cada método lê uma instância via entrada padrão (que por sua vez está estruturada de acordo com a descrição do trabalho), resolve o problema e registrar as métricas de execução.

\subsection{Geração das instâncias}
Para a geração dos grafos foi utilizada a biblioteca \texttt{jngen}~\footnote{\url{https://github.com/ifsmirnov/jngen}}, uma biblioteca feita para gerar testes para programação competitiva. Ela permite gerar grafos com número de vértices e arestas definidos, além de garantir conectividade (todas as instânicas geradas são conexas). Um programa foi criado apenas para gerar as instâncias, que por sua vez é executado antes dos solucionadores.

Foram considerados os seguintes número de vértices nas instâncias: $n \in \{5,10,15,$ $20,25\}$. Para cada valor de $n$, foram gerados grafos com três faixas de densidade: esparsa, média e densa, definidas em de intervalos de número de arestas. Dentro de cada faixa, o número de arestas é escolhido aleatoriamente dentro do intervalo correspondente. A Tabela \ref{tab:faixas-de-densidade} mostra as faixas de densidade para cada valor de $n$.

\begin{table}[H]
    \centering
    \begin{tabular}{|c|c|c|c|}
        \hline
        $n$ & Esparso & Médio & Denso \\
        \hline
        5  & $4$--$14$  & $15$--$4$  & $5$--$10$     \\
        10 & $9$--$29$  & $30$--$22$ & $23$--$45$    \\
        15 & $14$--$44$ & $45$--$52$ & $53$--$105$   \\
        20 & $19$--$59$ & $60$--$94$ & $95$--$190$   \\
        25 & $24$--$74$ & $75$--$149$ & $150$--$300$ \\
        \hline
    \end{tabular}
    \caption{Intervalos de número de arestas por faixa de densidade para cada valor de $n$.}
    \label{tab:faixas-de-densidade}
\end{table}

As faixas de densidade foram definidas com o número de arestas $m$ em função do número de vértices $n$. Seja $m_{\max} = \frac{n(n-1)}{2}$ o número máximo de arestas em um grafo simples e não direcionado. As três faixas foram definidas como:

\begin{itemize}
    \item \textbf{Esparso}: $n-1 \leq m < 3n$;
    \item \textbf{Médio}: $3n \leq m < \frac{n(n-1)}{4}$;
    \item \textbf{Denso}: $\frac{n(n-1)}{4} \leq m \leq m_{\max}$.
\end{itemize}

O limite inferior $n-1$ para grafos esparsos garante conectividade geração. $3n$ e $\frac{n(n-1)}{4}$  foram escolhidos arbitrariamente como limites inferior e superior, respectivamente, para determinar grafos médios, enquanto $m_{\max}$, por ser o número máximo de arestas, foi definido como limite superior dos grafos densos. 

É válido salientar que as fórmulas para $m$ em função de $n$ não são consistentes (ex: para $n=5$, $3n>\frac{n(n-1)}{4}$), então ajustes manuais foram necessários para chegar ao resultado final descrito na Tabela~\ref{tab:faixas-de-densidade}. Dentro de cada faixa, o número de arestas é escolhido aleatoriamente dentro do intervalo correspondente para cada instância gerada. Para cada par $(n,\text{densidade})$ foram geradas três instâncias independentes, totalizando $5 \cdot 3 \cdot 3 = 45$ grafos. Cada instância foi salva em um arquivo e utilizada como entrada para ambos os solucionadores.

\subsection{Métricas}
As métricas analisadas foram:

\begin{itemize}
    \item tempo de execução (em milissegundos);
    \item uso máximo de memória durante a execução do algoritmo (em kilobytes).
\end{itemize}

Cada solucionador foi responsável por mensurar suas próprias métricas, que por sua vez começaram a ser medidas logo após a leitura do grafo pela entrada padrão e finalizaram logo após a solução final ser impressa. O objetivo é desconsiderar tempo e memória gastas para inicialização do programa e leitura do grafo. O tempo de execução foi medido usando meios fornecidos pela biblioteca padrão do sistema para obter os \textit{timestamps} do início do final, com o resultado sendo calculado, em milissegundos, pela diferença entre ambos.

Já para a memória, foi adotada a seguinte estratégia: inicialmente, mede-se a memória utilizada atualmente pelo processo; após a execução do algoritmo, obtém-se o pico de memória atingido pelo processo. A diferença entre esses dois valores corresponde à quantidade máxima adicional de memória utilizada pelo algoritmo durante sua execução. Para cada faixa $(n,\text{densidade})$, o valor reportado corresponde à média aritmética das três instâncias geradas.

\subsection{Execução dos experimentos e Ambiente}
Cada execução de um algoritmo sobre uma instância foi realizada em um processo separado, garantindo que as métricas não fossem influenciadas por execuções anteriores. Não foi imposto limite de tempo para as execuções, pois os tamanhos considerados são computacionalmente viáveis para ambos os métodos.

Os experimentos foram realizados em uma máquina com as seguintes características:
\begin{itemize}
    \item CPU: AMD $\text{Ryzen}^{\circledR}$ 5 6600H (3,3 GHz);
    \item Memória: 16 GB DDR5 (4800 MT/s);
    \item Sistema operacional: $\text{Windows}^{\circledR}$ 11;
    \item Linguagem utilizada: C++
    \item Compilador: \texttt{g++} versão 14.2.0;
    \item Flags de compilação: \texttt{-O2}.
\end{itemize}
Os dados brutos dos experimentos, bem como o código utilizado para a geração das instâncias e para a execução dos algoritmos, além das bibliotecas externas, estão disponíveis no repositório do projeto (\href{https://github.com/izak-ag/IC0004-Trabalho-Pratico}{https://github.com/izak-ag/IC0004-Trabalho-Pratico}). O projeto foi construído para também ser compilável e executável em Linux.

    \section{Resultados} \label{sec:resultados}

Nesta seção são apresentados e discutidos os resultados obtidos para os dois métodos avaliados, considerando separadamente os três cenários de densidade dos grafos: esparso, médio e denso. Para cada cenário são analisados o tempo de execução e o uso de memória em função do número de vértices $n$. Para facilitar a visualização das diferenças de ordem de grandeza, todos os gráficos foram plotados com escala logarítmica no eixo correspondente ao tempo/memória (eixo y).

\subsection{Cenário esparso}

A \reffig{fig:esparso-tempo} apresenta o tempo médio de execução dos dois algoritmos em grafos esparsos. Observa-se que a abordagem baseada em programação dinâmica (BMK) apresenta tempos desprezíveis para $n \leq 10$ e cresce de forma acentuada a partir de $n = 15$, alcançando aproximadamente 12 segundos para $n = 25$. Por outro lado, a abordagem baseada em redução para SAT (RedSAT) apresenta crescimento muito mais rápido, passando de dezenas de milissegundos para centenas de segundos quando $n$ aumenta de 15 para 25.

A \reffig{fig:esparso-memoria} mostra o consumo de memória adicional dos dois algoritmos em grafos esparsos. O uso de memória do BMK cresce de forma moderada com $n$, enquanto o RedSAT apresenta um consumo inicial significativamente maior, mas com crescimento mais suave. Para grafos maiores, os valores são semelhantes.

\begin{figure}[H]
    \centering
    \grafico{sparse}{time}{Tempo de execução (ms)}
    \caption{Tempo médio de execução em função de $n$ para grafos esparsos.}
    \label{fig:esparso-tempo}
\end{figure}

\begin{figure}[H]
    \centering
    \grafico{sparse}{memory}{Memória adicional (KB)}[south east]
    \caption{Uso médio de memória adicional em função de $n$ para grafos esparsos.}
    \label{fig:esparso-memoria}
\end{figure}

\subsection{Cenário médio}

A \reffig{fig:medio-tempo} apresenta os tempos de execução para grafos de densidade média. O comportamento observado é semelhante ao do cenário esparso. O BMK mantém desempenho superior em toda a faixa de valores testados, enquanto o RedSAT sofre uma explosão combinatória significativa para $n \geq 20$.

Em relação à memória, conforme mostrado na \reffig{fig:medio-memoria}, BMK apresenta maior sensibilidade ao aumento de $n$ do que o RedSAT, especialmente para valores maiores de $n$, embora ambos permaneçam dentro de limites modestos em termos absolutos.

\begin{figure}[H]
    \centering
    \grafico{medium}{time}{Tempo de execução (ms)}
    \caption{Tempo médio de execução em função de $n$ para grafos de densidade média.}
    \label{fig:medio-tempo}
\end{figure}

\begin{figure}[H]
    \centering
    \grafico{medium}{memory}{Memória adicional (KB)}[south east]
    \caption{Uso médio de memória adicional em função de $n$ para grafos de densidade média.}
    \label{fig:medio-memoria}
\end{figure}

\subsection{Cenário denso}

A \reffig{fig:denso-tempo} mostra os tempos de execução para grafos densos. Assim como nos outros cenários, o BMK supera amplamente o RedSAT em todos os tamanhos considerados. O aumento da densidade não altera significativamente o comportamento assintótico dos tempos de execução, indicando que o fator dominante é o número de vértices.

A \reffig{fig:denso-memoria} apresenta o uso de memória adicional. Neste cenário, o BMK apresenta crescimento mais acentuado do consumo de memória, chegando a ultrapassar o RedSAT para valores maiores de $n$. Isso está associado ao maior número de estados viáveis quando o grafo é mais denso.

\begin{figure}[H]
    \centering
    \grafico{dense}{time}{Tempo de execução (ms)}
    \caption{Tempo médio de execução em função de $n$ para grafos densos.}
    \label{fig:denso-tempo}
\end{figure}

\begin{figure}[H]
    \centering
    \grafico{dense}{memory}{Memória adicional (KB)}[south east]
    \caption{Uso médio de memória adicional em função de $n$ para grafos densos.}
    \label{fig:denso-memoria}
\end{figure}

\subsection{Discussão geral}

De forma geral, os resultados mostram que BMK é significativamente mais eficiente em tempo de execução para todos os tamanhos e densidades considerados, enquanto o RedSAT apresenta crescimento muito mais rápido. Em termos de memória, ambos os métodos apresentam consumo relativamente baixo em valores absolutos para os tamanhos considerados neste trabalho, não ultrapassando alguns megabytes mesmo para as maiores instâncias. 

No entanto, o uso de memória é alto para valores $n$ maiores do que utilizados neste trabalho, pois ambas abordagens possuem complexidade espacial exponencial. No caso do BMK, o consumo espacial explode, uma vez que o algoritmo armazena uma tabela de tamanho $n\cdot2^n$. Por exemplo, para $n=30$ tem-se $30\cdot2^{30} > 32$ GB mesmo assumindo apenas 1 byte por entrada, tornando essa abordagem inviável em termos de memória para valores maiores de $n$.

Já o RedSAT, como observado nos experimentos, tende a apresentar crescimento mais suave do consumo de memória, embora também possa se tornar inviável para instâncias suficientemente grandes dependendo do tamanho da fórmula gerada. O BMK é mais sensível ao número de vértices, enquanto o RedSAT apresenta comportamento mais estável em relação a essa variável.

O principal ponto fraco do RedSAT em relação ao desempenho temporal é o custo da própria redução. A formulação gera $O(n^2)$ variáveis e $O(n^3)$ cláusulas, produzindo instâncias de SAT grandes, que são difíceis de resolver mesmo para os mais modernos solucionadores, inclusive o CDCL. De um modo geral, os resultados indicam que a redução para SAT não se mostra competitiva em termos de desempenho prático se comparado a um algoritmo especializado, ao menos no intervalo de tamanhos considerado neste trabalho.
    \section{Conclusão}

Neste trabalho foi estudado o Problema do Caminho Hamiltoniano de forma teórica e experimental. Inicialmente, foi apresentada a definição formal do problema e provada formalmente o seu pertencimento à classe NP e sua NP-completude por meio de reduções, começando a partir do 3-SAT para o problema do Caminho Hamiltoniano Dirigido, que por sua vez foi reduzido para o caso não direcionado, que é o foco deste trabalho. 

Também foram discutidos, de forma mais introdutória, métodos clássicos como backtracking e branch-and-bound, bem como heurísticas (como métodos gulosos, método de Pósa e metaheurísticas) e aproximações propostas na literatura, que embora mais rápidas não garantem exatidão. Essas abordagens foram apresentadas principalmente para contextualização. Em seguida, foram analisadas duas abordagens exatas para sua resolução: um algoritmo baseado em programação dinâmica (BMK) e uma abordagem baseada em redução para SAT resolvida por um solucionador CDCL (RedSAT).

Os resultados dos experimentos mostraram que o  BMK apresenta desempenho superior em termos de tempo de execução para todas as instâncias consideradas. Por outro lado, a abordagem baseada em SAT apresentou crescimento significativamente mais rápido do tempo de execução.

Em termos de uso de memória, ambas as abordagens apresentaram consumo moderado para os tamanhos considerados no experimento. Entretanto, a análise teórica indica que o consumo de memória do BMK explode a partir de valores de $n$ um pouco maiores do que os deste trabalho. Já o uso de memória no RedSAT tende a crescer de forma mais suave em relação ao BMK, ainda que o consumo seja relativamente bem alto para valores de $n$ menores.

Os resultados indicam que, para instâncias pequenas e médias do Problema do Caminho Hamiltoniano, um algoritmo exato especializado é preferível a uma formulação genérica via SAT, tanto por razões de eficiência quanto de simplicidade de implementação. No entanto, para valores de $n$ suficientemente grandes ambas abordagens são computacionalmente inviáveis devido às suas complexidades exponenciais.

Como trabalhos futuros, podem ser analisados o uso de heurísticas e algoritmos aproximados para instâncias maiores, onde também poem ser analisada a precisão dos mesmos (afinal, verificar a corretude da resposta é polinomial). Também podem ser investigados métodos híbridos, que combinem métodos exatos (como programação dinâmica) com heurísticas e/ou aproximações.
    
    \bibliographystyle{unsrt}
    \bibliography{referencias}
\end{document}